\documentclass{article}%
\usepackage[T1]{fontenc}%
\usepackage[utf8]{inputenc}%
\usepackage{lmodern}%
\usepackage{textcomp}%
\usepackage{lastpage}%
\usepackage{geometry}%
\geometry{margin=2cm}%
\usepackage{expex}%
\usepackage[T1]{fontenc}%
\usepackage[utf8]{inputenc}%
\usepackage[spanish]{babel}%
%
\title{Análisis Lingüístico: Gramatica Normativa Kaqchikel}%
\author{Generado por el TFM de Elías Carrasco}%
%
\begin{document}%
\normalsize%
\maketitle%
\section{Page 205}%
\label{sec:Page205}%
\par\medskip%
\ex
\textit{Lo quiere} \\
\begingl
\gla INC-B3 //
\glb INC-B3 //
\glc \textbf{A3-querer-SC} //
\endgl
\xe
%
\medskip%
\subsection{Negación de una acción}%
\label{subsec:Negacindeunaaccin}%

%
\par\medskip%
\ex
\textit{No lo quiere} \\
\begingl
\gla k'am //
\glb k'am chi-ø y-oche-j //
\glc \textbf{NEG} \textbf{INC-B3} \textbf{A3-querer-SC} //
\endgl
\xe
%
\medskip%
\subsection{Oración base (acción en potencial)}%
\label{subsec:Oracinbase(accinenpotencial)}%

%
\par\medskip%
\ex
\textit{Él se bañará} \\
\begingl
\gla hoq ø achn-oq //
\glb hoq ø achn-oq //
\glc \textbf{POT-B3} \textbf{bañar-SC} \textbf{naq} //
\glft \textbf{Sintaxis:} [hoq ø achn-oq]\textsubscript{PRO:él} //
\endgl
\xe
%
\medskip%
\subsection{Negación de una acción}%
\label{subsec:Negacindeunaaccin}%

%
\par\medskip%
\ex
\textit{No se bañará} \\
\begingl
\gla man hoq ø achn-oq //
\glb man hoq ø achn-oq //
\glc \textbf{no} \textbf{POT-B3} \textbf{bañar-SC} //
\glft \textbf{Sintaxis:} [naq]\textsubscript{PRO:él} //
\endgl
\xe
%
\medskip%
\subsection{Oración base (acción en incompletivo)}%
\label{subsec:Oracinbase(accinenincompletivo)}%

%
\par\medskip%
\ex
\textit{Me iré mañana} \\
\begingl
\gla ch-in toj //
\glb ch-in toj yekal //
\glc \textbf{INC-B1SG} \textbf{ir} \textbf{mañana} //
\endgl
\xe
%
\medskip%
\subsection{Negación de una acción}%
\label{subsec:Negacindeunaaccin}%

%
\par\medskip%
\ex
\textit{No me iré mañana} \\
\begingl
\gla maj hin toq //
\glb maj hin toq yekal //
\glc \textbf{no} \textbf{A1SG} \textbf{ir} \textbf{mañana} //
\endgl
\xe
%
\medskip%
\subsection{Oración base (predicado no verbal)}%
\label{subsec:Oracinbase(predicadonoverbal)}%

%
\par\medskip%
\ex
\textit{Ella está parada} \\
\begingl
\gla lek-an Ø //
\glb lek-an Ø aj //
\glc \textbf{parado-POS} \textbf{B3} \textbf{DIR:hacia arriba} //
\glft \textbf{Sintaxis:} [ix]\textsubscript{PRO:ella} //
\endgl
\xe
%
\medskip%
\subsection{Negación de un predicado no verbal}%
\label{subsec:Negacindeunpredicadonoverbal}%

%
\par\medskip%
\ex
\textit{} \\
\begingl
\gla man lekan-oq //
\glb man Ø //
\glc \textbf{lekan-oq} \textbf{aj ix} //
\endgl
\xe
%
\medskip

%
\end{document}